\documentclass[10pt,a4paper]{amsart}
\usepackage[margin=19mm]{geometry}
\usepackage{amsmath,amssymb,mathtools}
\usepackage[scr=rsfs,cal=euler]{mathalfa}
\usepackage{xfrac}
\usepackage[unicode,hidelinks,bookmarks=false]{hyperref}
\newcommand{\ie}{i.e.\ }
\newcommand{\cf}{cf.\ }
\newcommand{\resp}{respectively\ }
\newcommand{\Subsection}{\subsection}
\providecommand{\nobreakdash}{\nobreak\hbox{-}}
\setlength{\emergencystretch}{2em}
\multlinegap0pt
\usepackage[shortlabels]{enumitem}
\usepackage{amssymb}
\usepackage{mathtools}
\newtheorem{lem}{Lemma}
\newcommand{\real}{\mathbb{R}}
\title{\vspace*{-1.5cm} Submultiplicative norms and filtrations on section rings}
\author{Siarhei Finski}
\date{Source: arXiv:2210.03039v3 (CC BY 4.0)}
\begin{document}
\maketitle

\noindent\emph{Source and licence.} \vspace*{-1.5cm} Submultiplicative norms and filtrations on section rings. Version arXiv:2210.03039v3 (CC BY 4.0), distributed by arXiv under the Creative Commons Attribution 4.0 International licence, http://creativecommons.org/licenses/by/4.0/, as recorded in the arXiv licence metadata for this exact version and shown on the abstract page https://arxiv.org/abs/2210.03039v3. Full licence notice: \texttt{licenses/arxiv-2210.03039-CC-BY-4.0.txt}.\par\smallskip

\noindent\textit{Editorial scope.} Counted unit: the article's lem lem\_log\_spec\_ray (source0.tex lines 877--888). The article's complete original proof of it is delivered with it (source0.tex lines 889--988, one \texttt{proof} environment of 100 lines). No mathematics is written, repaired or re-derived: the statement and the proof are the article's own lines, in the article's own notation, and no retained line is substituted or re-flowed.  Retained uncounted as the source-local context the proof needs: lines 769--774; lines 776--778; lines 791--793; lines 864--876.  Nothing was omitted from any retained span, so the package carries no omission record.  No retained line contains a citation, so the wrapper carries no bibliography and no bibliography entry.  No retained line contains a figure environment, an image-inclusion command or drawing code, so no image is delivered and the package declares no asset dependency.  Editorial formatting only, in addition: an AMS \texttt{amsart} preamble that restates the article's own declarations and macros used by the retained lines, namely the environments \texttt{lem} on separate counters, together with the standard packages those lines need.  The retained environments print in this document as Lemma 1 in order of appearance, whereas the article prints the same items with the numbering of its own full document.  The complete native source of the article as distributed in the arXiv e-print for this exact version is bundled unchanged as \texttt{sources/131716/source0.tex}.
	\begin{equation}\label{eq_log_rel_spec}
		\lambda_j
		:=
		\sup_{\substack{W \subset V \\ \dim W = j}} 
		\inf_{w \in W \setminus \{0\}} \log \frac{\| w \|_2}{\| w \|_1}.
	\end{equation}

	\begin{equation}\label{eq_log_rel_spec02}
		\lambda_j(N_1, N_2) + \lambda_{\dim V}(N_2, N_3) \leq \lambda_j(N_1, N_3) \leq \lambda_j(N_1, N_2) + \lambda_1(N_2, N_3),
	\end{equation}

	\begin{equation}\label{eq_john_ellips}
		N_V^H \leq N_V \leq \sqrt{\dim V} \cdot N_V^H.
	\end{equation}

	\begin{equation}\label{eq_ray_norm_defn0}
		\| f \|_{V, \mathcal{F}}^t :=
		\inf 
		\Big\{
			\sum
			e^{- t \mu_i} \cdot 
			\| f_i \|_V
			\,
			:
			\,
			f = \sum f_i, \, f_i \in \mathcal{F}^{\mu_i} V
		\Big\}.
	\end{equation}

	\begin{lem}\label{lem_log_spec_ray}
		For any $t \geq s \geq 0$ and $i = 1, \ldots, \dim V$, the following estimate holds
		\begin{equation}
			\Big| \lambda_i(N_{V, \mathcal{F}}^t, N_{V, \mathcal{F}}^s) - (t - s) \cdot e_{\mathcal{F}}(i) \Big|
			\leq
			\log \dim V,
		\end{equation}
		where $e_{\mathcal{F}}(i)$ are the jumping numbers of the filtration $\mathcal{F}$, defined as
		\begin{equation}\label{eq_defn_jump_numborig}
			e_{\mathcal{F}}(i) := \sup \Big\{ t \in \real : \dim \mathcal{F}^t H^0(X, L^{\otimes k}) \geq i \Big\}.
		\end{equation}
	\end{lem}

	\begin{proof}
		Let us first show that for any $t \geq s \geq 0$, $f \in V$ the following identity holds
		\begin{equation}\label{eq_ray_norm_defn2}
			\| f \|_{V, \mathcal{F}}^t =
			\inf 
			\Big\{
				\sum
				e^{- (t - s) \mu_i} \cdot 
				\| f_i \|_{V, \mathcal{F}}^s
				\,
				:
				\,
				f = \sum f_i, \, f_i \in \mathcal{F}^{\mu_i} V
			\Big\}.
		\end{equation}
		Clearly, once (\ref{eq_ray_norm_defn2}) is established, it would be enough to prove Lemma \ref{lem_log_spec_ray} for $s = 0$.
		\par 
		First of all, the inequality 
		\begin{equation}\label{eq_ray_norm_defn2aa1}
			\| f \|_{V, \mathcal{F}}^t \geq
			\inf 
			\Big\{
				\sum
				e^{- (t - s) \mu_i} \cdot 
				\| f_i \|_{V, \mathcal{F}}^s
				\,
				:
				\,
				f = \sum f_i, \, f_i \in \mathcal{F}^{\mu_i} V
			\Big\},
		\end{equation}
		follows directly from the bound $\| f_i \|_{V, \mathcal{F}}^s \leq e^{- s \mu_i} \cdot \| f_i \|_V$ and (\ref{eq_ray_norm_defn0}).
		\par 
		To establish the inverse inequality, we first remark that it is clear (by the use of projection operator) that if $f  \in \mathcal{F}^{\mu} V$, $\mu \in \real$, then in (\ref{eq_ray_norm_defn0}), it is sufficient to consider the decompositions with $f_i \in \mathcal{F}^{\mu_i} V$, $\mu_i \geq \mu$.
		For any $\epsilon > 0$, we now consider a decomposition $f = \sum f_i^{\epsilon}$, $f_i^{\epsilon} \in \mathcal{F}^{\mu_i} V$, verifying
		\begin{equation}
			\sum_{i}
				e^{- (t - s) \mu_i} \cdot 
				\| f_i^{\epsilon} \|_{V, \mathcal{F}}^s
			\leq
			\inf 
			\Big\{
				\sum_{i}
				e^{- (t - s) \mu_i} \cdot 
				\| f_i \|_{V, \mathcal{F}}^s
				\,
				:
				\,
				f = \sum f_i, \, f_i \in \mathcal{F}^{\mu_i} V
			\Big\}
			+
			\epsilon.
		\end{equation}
		We also consider decompositions $f_i^{\epsilon} = \sum f_{i, j}^{\epsilon}$, $f_{i, j}^{\epsilon} \in \mathcal{F}^{\mu_{i, j}} V$, $\mu_{i, j} \geq \mu_i$, verifying 
		\begin{equation}
			\sum_{j}
				e^{- s \mu_{i, j}} \cdot 
				\| f_{i, j}^{\epsilon} \|_V
			\leq
			\| f_i^{\epsilon} \|_{V, \mathcal{F}}^s
			+
			\epsilon.
		\end{equation}
		Clearly, since $\mu_{i, j} \geq \mu_i$, we then have
		\begin{equation}
			\sum_{i, j}
				e^{- (t - s) \mu_i} 
				\cdot
				e^{- s \mu_{i, j}} \cdot 
				\| f_{i, j}^{\epsilon} \|_V
			\geq
			\sum_{i, j}
				e^{- t  \mu_{i, j}} 
				\cdot 
				\| f_{i, j}^{\epsilon} \|_V
			\geq 
			\| f \|_{V, \mathcal{F}}^t,
		\end{equation}
		which, along with the above estimates, imply the inverse inequality to (\ref{eq_ray_norm_defn2aa1}), when $\epsilon \to 0$.
		\par 
		Let assume first that $N_V$ is Hermitian. We then show the following stronger statement
		\begin{equation}\label{eq_spec_herm}
			\lambda_i(N_{V, \mathcal{F}}^t, N_V) = t \cdot e_{\mathcal{F}}(i).
		\end{equation}
		For simplicity of the presentation, we assume that the numbers $e_{\mathcal{F}}(i)$ are different.
		Remark that if $x \in V$ is orthogonal (with respect to the scalar product associated with $N_V$) to $\mathcal{F}^{e_{\mathcal{F}}(i - 1)} V$, then $\| x \|_{V, \mathcal{F}}^t \geq \| x \|_V \cdot \exp(- t e_{\mathcal{F}}(i))$.
		In particular, since in any $i$-dimensional subspace of $V$, there is an element orthogonal to the $i - 1$-dimensional subspace $\mathcal{F}^{e_{\mathcal{F}}(i - 1)} V$, we have $\lambda_i(N_{V, \mathcal{F}}^t, N_V) \leq t  \cdot e_{\mathcal{F}}(i)$.
		However, by taking a subspace $W := \mathcal{F}^{e_{\mathcal{F}}(i)} V$ in (\ref{eq_log_rel_spec}), we obtain $\lambda_i(N_{V, \mathcal{F}}^t, N_V) \geq t \cdot e_{\mathcal{F}}(i)$.
		A combination of the two estimates imply (\ref{eq_spec_herm}).
		\par 
		Now, for general norms $N_V$, we consider the norm $N_V^H$ as in (\ref{eq_john_ellips}).
		Let us denote by $N_{V, \mathcal{F}}^{H, t}$ the ray of norms emanating from $N_V^H$ as in (\ref{eq_ray_norm_defn0}).
		Clearly, by (\ref{eq_john_ellips}), for any $t \in [0, +\infty[$, we have 
		\begin{equation}
			d_{+ \infty}(N_{V, \mathcal{F}}^{H, t}, N_{V, \mathcal{F}}^t) \leq \frac{1}{2} \log \dim V,
			\qquad
			d_{+ \infty}(N_V, N_V^H) \leq \frac{1}{2} \log \dim V.
		\end{equation}
		By (\ref{eq_log_rel_spec02}) and (\ref{eq_spec_herm}), we conclude the proof of Lemma \ref{lem_log_spec_ray}.
	\end{proof}

\end{document}
